\section{Glossario}

In questa sezione sono raccolti i termini che potrebbero risultare meno immediati per il lettore del manuale.

\term{\hypertarget{glossario:amministratore}{Amministratore}}
Tipo di utente con permessi dedicati alla consultazione e al monitoraggio degli ordini, senza accesso alle funzionalit\`a di creazione degli ordini.

\term{\hypertarget{glossario:api-key}{API Key}}
Codice segreto usato per autorizzare un'applicazione a comunicare con un servizio esterno. In SmartOrder serve per collegare il sistema ai servizi OpenAI.

\term{\hypertarget{glossario:autenticazione}{Autenticazione}}
Processo con cui il sistema verifica l'identit\`a dell'utente tramite credenziali, ad esempio username e password.

\term{\hypertarget{glossario:browser}{Browser}}
Programma usato per navigare sul web e aprire applicazioni accessibili tramite pagina web, ad esempio Google Chrome, Microsoft Edge, Mozilla Firefox o Safari.

\term{\hypertarget{glossario:chat}{Chat}}
Area principale dell'applicazione in cui l'utente scrive messaggi o invia input vocali per costruire e modificare il proprio ordine.

\term{\hypertarget{glossario:chatbot}{Chatbot}}
Sistema software progettato per interagire con l'utente tramite messaggi, simulando una conversazione naturale.

\term{\hypertarget{glossario:container}{Container}}
Processo avviato da Docker che esegue una parte di SmartOrder in un ambiente isolato. 
Durante l’installazione, l’applicazione funziona solo se i container risultano \textit{in esecuzione}.

\term{\hypertarget{glossario:docker}{Docker}}
Piattaforma software che permette di creare, avviare e gestire container.

\term{\hypertarget{glossario:file-env}{File \texttt{.env}}}
File di configurazione con variabili d’ambiente usate dall’applicazione (es. chiavi API e password).
Contiene dati \textit{sensibili} e non deve essere condiviso.

\term{\hypertarget{glossario:git}{Git}}
Programma usato per scaricare (\vr{clonare}) un progetto da GitHub tramite comandi da terminale.

\term{\hypertarget{glossario:llm}{LLM}}
Sigla di \textit{Large Language Model}. Indica un modello di intelligenza artificiale addestrato a comprendere e generare linguaggio naturale.

\term{\hypertarget{glossario:login}{Login}}
Operazione di accesso al sistema tramite inserimento delle credenziali personali.

\term{\hypertarget{glossario:logout}{Logout}}
Operazione con cui l'utente esce dall'applicazione e termina la propria sessione di accesso.

\term{\hypertarget{glossario:openai}{OpenAI}}
Azienda che fornisce modelli e servizi di intelligenza artificiale utilizzati da SmartOrder per la gestione della chat.

\term{\hypertarget{glossario:paginazione}{Paginazione}}
Suddivisione di una lista lunga in pi\`u pagine, per facilitarne la consultazione.

\term{\hypertarget{glossario:readme}{README}}
File testuale di un progetto che contiene in genere informazioni introduttive, istruzioni di utilizzo e indicazioni sulla struttura del software.

\term{\hypertarget{glossario:repository}{Repository}}
Pagina del progetto su GitHub che contiene i file: da lì puoi scaricare lo \texttt{.zip} o clonare con Git.

\term{\hypertarget{glossario:sessione}{Sessione}}
Periodo di tempo in cui un utente rimane autenticato all'interno del sistema dopo aver effettuato il login.